\hmsubsection{Robotics Applications}{UMA}{OAH}\label{sec:robotics-applications}

Robotic manipulation is widely used in applications where objects must be
handled with repeatable motion and limited human intervention. Typical examples
include industrial pick-and-place cells, educational robot arms, laboratory
automation, and small-scale sorting demonstrators. These systems usually combine
a perception module that identifies objects in the workspace with a motion
module that converts object positions into actuator commands.

The Autonomia system follows this general structure, but at a scale suitable for
a bachelor project. The task is constrained to coloured spherical objects, a
fixed workspace, and a limited number of destination bins. These constraints make
it possible to use deterministic computer vision and closed-form inverse
kinematics rather than more complex industrial approaches such as high-degree-of-
freedom motion planning or learning-based grasp planning. The related robotics
literature therefore supports the central design choice of treating the system as
a controlled pick-and-place application, where reliability and explainability are
prioritised over general-purpose manipulation.
